\chapter{Randomisierte Algorithmen}
Wir wollen den Begriff des \textit{deterministischen Algorithmus} einführen. Wikipedia beschreibt diesen Begriff (treffend) wie folgt:

\textit{Ein deterministischer Algorithmus ist ein Algorithmus, bei dem nur definierte und reproduzierbare Zustände auftreten. Für die gleiche Eingabe folgt auch immer die gleiche Ausgabe und zusätzlich wird die gleiche Folge an Zuständen durchlaufen. Zu jedem Zeitpunkt ist der nachfolgende Abarbeitungsschritt des Algorithmus eindeutig festgelegt. Das bedeutet auch, dass alle Zwischenergebnisse innerhalb des Algorithmus immer gleich sind.}

Die Algorithmen, die wir bislang betrachtet haben, waren deterministisch. Sie haben auf gleichen Eingaben immer genau gleich gearbeitet und dabei stets die exakt gleiche Folge von Zuständen durchlaufen. Der nächste Schritt des Programms war dabei immer (von Anfang an) eindeutig festgelegt.

Ein randomisierter Algorithmus wollen wir uns vorstellen, wie ein deterministischer Algorithmus, der mindestens eine Zufällig gewählte Zahl erhält, welche dann das Resultat der Berechnung beeinflusst. Man kann sich auch vorstellen, dass der Algorithmus an mindestens einer Stelle eine Münze wirft und seine Berechnung / Entscheidung (z.B. bei einem \mintinline{python}{if}) vom Ausgang dieses zufälligen Münzwurfs macht.

Ziel dieses Kapitels ist zu zeigen, dass randomisierte Algorithmen deutlich effizienter als bestmögliche deterministische Algorithmen sein können.

Dazu werden wir uns im Detail mit einem Beispiel eines randomisierten Algorithmus beschäftigen.

Das Beispiel und die Notation wurde aus dem exzellenten Buch \textit{Theoretische Informatik, 5. Auflage, Abschnitt 8.3} von Juraj Hromkovi\v{c} übernommen.

\section{Ein randomisiertes Kommunikationsprotokoll}
Sei $R_1$ ein Rechner, der sich an der ETH in Zürich befindet und sei $R_2$ ein Rechner am MIT in Boston, USA. Ursprünglich besassen diese Rechner je eine Datenbank mit exakt gleichem Inhalt. Mit der Zeit hat sich der Inhalt geändert, aber man hat sich bemüht, jeweils die gleichen Änderungen auf beiden Rechnern durchzuführen, um idealerweise die gleiche Datenbank auf beiden Rechnern zu erhalten. Nach einiger Zeit wollen wir prüfen, ob $R_1$ und $R_2$ tatsächlich noch dieselben Daten haben. Sei $n$ die Grösse der Datenbank in Bits. Wir wollen hier konkret eine Grösse von $n = 10^{16}$ betrachten. Unser Ziel ist es, einen Kommunikationsalgorithmus (ein Protokoll) zu
entwerfen, der feststellt, ob sich die Inhalte der Datenbanken von $R_1$ und $R_2$ unterscheiden.

Das offensichtliche Vorgehen, um diese Frage zu beantworten, ist z.B. den Datensatz in Zürich Bit für Bit nach Boston zu senden, damit die Leute am MIT ihren Datensatz Bit für Bit mit dem Datensatz von Zürich vergleichen können. Dieses Vorgehen benötigt einen Austausch von genau $n$ Bits. Man kann beweisen, dass es keinen deterministischen Algorithmus gibt, der mit einem Austausch von weniger als $n$ Bit auskommt (bereits $n-1$ Bits reichen nicht aus). Offensichtlich ist diese Situation nicht sehr befriedigend, da es nicht praktisch ist, solch grosse Datenmengen zu senden. Desweiteren ist die Wahrscheinlichkeit, dass bei der Übertragung solch vieler Bits mindestens ein Bit falsch übertragen werden könnte, nicht zu vernachlässigen.

Im Folgenden wollen wir ein randomisiertes Kommunkationsprotokoll entwerfen, welches dieses Problem mit dem Austausch von wesentlich weniger Bits lösen kann. Dazu benötigen wir allerdings ein paar mathematische Tools, welche wir hier vorstellen wollen.

\subsection{Mathematische Tools}
In diesem Unterabschnitt stellen wir die nötigen mathematischen Werkzeuge vor, welche benötigt werden um das randomisierte Protokoll vollständig verstehen zu können.

Es ist wichtig, diese Tools gut zu verstehen.

\subsubsection{Fundamentalsatz der Arithmetik / Primfaktorzerlegung}
Jede natürliche Zahl lässt sich als Produkt von endlich vielen Primzahlen darstellen. Ordnet man diese Primzahlen der Grösse nach, so ist diese Darstellung sogar eindeutig.

Insbesondere kann man jede natürliche Zahl $n>0$ eindeutig darstellen als
\begin{align}
    n = p_{1}^{i_1}p_{2}^{i_2}\cdot ...\cdot p_{k}^{i_k},
\end{align}
wobei $p_1<p_2<...<p_k$ Primzahlen und $i_1, i_2, ..., i_k$ positive ganze Zahlen sind.

Beispiele:
\begin{align*}
    36 &= 2^2\cdot 3^2, \quad (p_1 = 2, p_2 = 3, i_1 = 2, i_2 = 2)\\
    99825 & = 3^1\cdot 5^2 \cdot 11^3, \quad (p_1 = 3, p_2 = 5, p_3 = 11, i_1 = 1, i_2 = 2, i_3 = 3)
\end{align*}

\begin{Exercise}%[difficulty={1}]

Finde die Primfaktorzerlegung der Zahl $1960$ und stelle sie in obiger Form dar.
\end{Exercise}

\begin{Answer}

\begin{align*}
    1960 = 2^3\cdot 5^1\cdot 7^2
\end{align*}
\end{Answer}


\subsubsection{Division mit Rest}
Zu zwei ganzen Zahlen $a$ und $b$, mit $b\neq 0$, existieren eindeutige ganze Zahlen $q$ und $r$, so dass
\begin{align}
    a = bq + r
\end{align}
und $0\leq r < |b|$.

Dabei ist $q$ das Resultat der ganzzahligen Division von $a/b$ und $r$ ist der Rest der Division.

Beispiel:

Seien $a = 17$ und $b=5$. Dann ist die ganzzahlige Division von $q=a/b=17/5=3$ mit einem Rest $r=a-bq=2$.

Daraus folgt direkt, dass wenn
\begin{align}
    a\; \text{mod}\; b = r,
\end{align}
dann existiert eine ganze Zahl $q$, so dass
\begin{align}
    a = bq + r.
\end{align}
Dabei kann $q$ durch die Berechnung $q = (a-r)/b$ ermittelt werden.

Beispiel:

Seien $a=79$ und $b=100$. Dann ist  $r = a\; \text{mod}\; b = 79\; \text{mod}\; 100 = 79$ und $q = (a-r)/b = (79-79)/100 = 0$.

\begin{Exercise}%[difficulty={1}]

\begin{enumerate}[label=(\alph*)]
    \item Bestimme die oben beschriebene Darstellung $a = bq + r$ für $a = 121$ und $b = 9$.
    \item Schreibe eine Python-Funktion \mintinline{python}{def division_mit_rest(a,b)}, welche für gegebene ganze Zahlen $a$ und $b$ mit $b \neq 0$ die Zahlen $q$ und $r$ berechnet und in einer Liste der Form \mintinline{python}{[q, r]} zurück gibt.
\end{enumerate}
\end{Exercise}

\begin{Answer}%[difficulty={1}]

\begin{enumerate}[label=(\alph*)]
    \item \begin{align*}
    121 = 9\cdot 13 + 4,
\end{align*}
wobei $a = 121$, $b = 9$, $q = 13$ und $r = 4$.
    \item \begin{minted}{python}
def division_mit_rest(a,b):
    r = a % b
    q = int((a-r)/b)
    return([q,r])
\end{minted}
\end{enumerate}
\end{Answer}


\subsubsection{$\text{Nummer}(x)$}
Sei $x=x_1x_2...x_n$, $x_i\in\left\{0,1\right\}$ für $i=1,2,...,n$ eine endliche Folge von binären Ziffern. Dann bezeichnet 
\begin{align}
    \text{Nummer}(x)=\displaystyle\sum_{i=1}^n x_i\cdot 2^{n-i}
\end{align}
die Darstellung von $x$ im Dezimalsystem.

Beispiele:
\begin{align*}
    x &= 00101 \Rightarrow \text{Nummer}(x) = 2^2 + 2^0 = 4 + 1 = 5\\
    x &= 01011 \Rightarrow \text{Nummer}(x) = 2^3 + 2^1  + 2^0 = 8+ 2 + 1 = 11
\end{align*}
\begin{Exercise}%[difficulty={1}]

\begin{enumerate}[label=(\alph*)]
    \item Berechne $\text{Nummer}(1011)$.
    \item Berechne $\text{Nummer}(110101)$.
\end{enumerate}
\end{Exercise}

\begin{Answer}
\begin{enumerate}[label=(\alph*)]
    \item \begin{align*}
    \text{Nummer}(1011) = 2^3 + 2^1 + 2^0 = 8 + 2 + 1 = 11
\end{align*}
    \item \begin{align*}
    \text{Nummer}(110101) = 2^5 + 2^4 + 2^2 + 2^0 = 32 + 16 + 4 + 1 = 53
\end{align*}
\end{enumerate}
\end{Answer}


\begin{Exercise}

Schreibe eine Python Funktion $\mintinline{python}{def Nummer(x)}$, welche $\text{Nummer}(x)$ für eine natürliche Zahl $x$ berechnet.
\end{Exercise}

\begin{Answer}
\begin{minted}[linenos]{python}
def Nummer(x):
    x = list(str(x))
    n = len(x)
    dec = 0
    for i, xi in enumerate(x):
        dec = dec + int(xi) * 2**(n-(i+1))
    return(dec)
\end{minted}
\end{Answer}



\subsubsection{Primzahlsatz und $\text{Prim}(n)$}
Euklid bewies bereits in der Antike, dass es unendlich viele Primzahlen gibt. Dies ist äquivalent zur Aussage, dass es keine grösste Primzahl gibt.

Wir wollen einmal die ersten $25$ natürlichen Zahlen auflisten und die Primzahlen unter ihnen unterstreichen
\begin{align*}
    0, 1, \underline{2}, \underline{3}, 4, \underline{5}, 6, \underline{7}, 8, 9, 10, \underline{11}, 12, \underline{13}, 14, 15, 16, \underline{17}, 18, \underline{19}, 20, 21, 22, \underline{23}, 24, ...
\end{align*}
Dank der Arbeit von Euklid wissen wir, dass in dieser Auflistung von natürlichen Zahlen immer wieder eine Primzahl auftauchen wird, egal wie gross die Zahlen in der Liste werden.

Im Folgenden wollen wir mit $\text{Prim}(n)$ die Anzahl der Primzahlen kleiner oder gleich $n$ bezeichnen. Beispielsweise sind die vier Primzahlen $2, 3, 5$ und $7$ alle Primzahlen, die kleiner oder gleich $9$ sind. Damit ist $\text{Prim}(9)$ = $4$.

Weitere Beispiele:

$\text{Prim}(10) = \text{Prim}(7) = \text{Prim}(9) = \text{Prim}(8) = 4$

$\text{Prim}(2) = 1$

$\text{Prim}(16) = 6$.


\begin{Exercise}
Bestimme die Zahl $\text{Prim}(33)$.
\end{Exercise}

\begin{Answer}

Die Primzahlen $\leq 33$ sind:

$2, 3, 5, 7, 11, 13, 17, 19, 23, 29, 31$.

Dies sind $11$ Zahlen. Damit ist $\text{Prim}(33) = 11$.
\end{Answer}



Eine natürliche Zahl $n\geq 2$ heisst Primzahl, wenn sie keinen echten Teiler hat. Deshalb könnte man intuitiv annehmen, dass grosse Primzahlen äusserst selten sind, da sehr viele Zahlen potentielle echte Teiler der Zahl zu sein scheinen.

Die Frage nach der Häufigkeit der Primzahlen ist eine alte Frage der Mathematik. Es hat sich herausgestellt, dass Primzahlen in einem gewissen Sinne überhaupt nicht selten sind, sondern sogar häufig auftauchen.

Ein starkes Resultat im Zusammenhang mit der Frage nach der Häufigkeit von Primzahlen ist der \textit{Primzahlsatz}. Diesen werden wir hier präsentieren, aber nicht beweisen. Der Beweis dieses Satzes ist sehr schwierig und wir werden ihn deshalb hier nicht führen.

\textbf{Primzahlsatz}

Der Primzahlsatz besagt
\begin{align}\label{eq:primzahlsatz}
    \lim_{n\to\infty} \frac{\text{Prim}(n)}{n/\ln (n)} = 1.
\end{align}
In Worten ausgedrückt, besagt der Satz, dass die Funktion $\text{Prim}(n)$ (Anzahl der Primzahlen, die kleiner oder gleich $n$ sind) für grosse $n$ (asymptotisch) gleich schnell wächst wie die Funktion $n/\ln (n)$.

Daraus folgt offensichtlich sofort, dass
\begin{align}\label{eq:quadrat}
    \lim_{n\to\infty} \frac{\text{Prim}(n^2)}{n^2/\ln (n^2)} = 1.
\end{align}
In der Form \ref{eq:quadrat} werden wir den Primzahlsatz später verwenden.
\begin{Exercise}
\begin{enumerate}[label=(\alph*)]
    \item Verwende den Primzahlsatz \ref{eq:primzahlsatz} um eine Approximation von $\text{Prim}(10^6)$ (die Anzahl der Primzahlen $\leq$ eine Million) zu finden.
    \item Bestimme selber den Wert $\text{Prim}(30)$ und vergleiche ihn mit der Schätzung, welche der Primzahlsatz \ref{eq:primzahlsatz} liefert. Warum ist es kein Problem, dass die Schätzung so schlecht zu sein scheint?
\end{enumerate}
\end{Exercise}:

\begin{Answer}
\begin{enumerate}[label=(\alph*)]
    \item Der Satz liefert eine Schätzung von $10^6 / \ln{(10^6)} \approx 72382.41$, also $\text{Prim}(10^6) = 72382$.
    \item $\text{Prim}(30) = 10$, da $2, 3, 5, 7, 11, 13, 17, 19, 23, 29$ die Primzahlen $\leq$ 30 sind.
Der Primzahlsatz \ref{eq:primzahlsatz} liefert die Schätzung $30 / \ln{(30)} \approx 8.82$.

Der Primzahlsatz macht nur eine Aussage über das asymptotische Verhalten von $\text{Prim}$. Die Abschätzung ist aber bereits für recht kleine Werte von $n$ relativ gut. Dass der Satz für solch kleine Werte wie $30$ keine ganz genaue Schätzung liefert ist kein Widerspruch.
\end{enumerate}
\end{Answer}

\subsubsection{Berechnen von  $\text{Prim}(n)$ in Python}
Der Algorithmus des \textit{Sieb des Eratosthenes} (englisch: Eratosthene's Sieve) ist bereits seit der Antike bekannt. Der Algorithmus bestimmt für eine gegebene natürliche Zahl $n$ alle Primzahlen, die kleiner oder gleich $n$. Man schreibt alle natürlichen Zahlen grösser gleich $2$ bis und mit $n$ auf:
\begin{align*}
    2, 3, 4, 5, 6, 7, 8, 9, 10, 11, 12, 13, 14, 15, 16, 17, 18, 19, 20, 21, 22, ..., n.
\end{align*}
Danach betrachtet man die erste Zahl in der Auflistung, also die Zahl $2$. Diese Zahl wurde noch nicht durchgestrichen, also ist sie eine Primzahl und das Sieb des Eratosthenes gibt diese Zahl aus. Danach wird jedes ganzzahlige Vielfache von $2$, dass kleiner oder gleich $n$ ist in der Liste gestrichen:
\begin{align*}
    2, 3, \cancel{4}, 5, \cancel{6}, 7, \cancel{8}, 9, \cancel{10}, 11, \cancel{12}, 13, \cancel{14}, 15, \cancel{16}, 17, \cancel{18}, 19, \cancel{20}, 21, \cancel{22}, ...
\end{align*}
Danach betrachtet der Algorithmus die nächste Zahl, also die Zahl $3$. Da diese noch nicht gestrichen wurde, ist $3$ eine Primzahl und wird ausgegeben. Nun werden alle ganzzahligen Vielfachen von $3$ ebenfalls gestrichen (falls sie nicht sowieso bereits gestrichen wurden):
\begin{align*}
    2, 3, \cancel{4}, 5, \cancel{6}, 7, \cancel{8}, \cancel{9}, \cancel{10}, 11, \cancel{12}, 13, \cancel{14}, \cancel{15}, \cancel{16}, 17, \cancel{18}, 19, \cancel{20}, \cancel{21}, \cancel{22}, ...
\end{align*}
Nun betrachtet der Algorithmus wieder die nächste Zahl in der Liste, also die Zahl $4$. Diese wurde aber bereits gestrichen und deshalb schaut sich der Algorithmus die nächste Zahl an (die Zahl $5$). Dieses Vorgehen wird wiederholt, bis der Algorithmus bei der Zahl $n$ ankommt.
\begin{Exercise}
\begin{enumerate}[label=(\alph*)]
\item Erkläre, warum dieser Algorithmus genau alle Primzahlen kleiner oder gleich $n$ ausgibt.
\item Wie kann das Sieb des Eratosthenes verwendet werden, um $\text{Prim}(n)$ zu bestimmen?
\end{enumerate}
\end{Exercise}


\begin{Answer}
\begin{enumerate}[label=(\alph*)]
\item Verlangt war nur eine informelle Erklärung. Wir zeigen hier etwas genauer:
\begin{center}
    eine natürliche Zahl $m$ ist eine Primzahl
    
    $\Longleftrightarrow$ 
    
    $m$ wird vom Sieb des Eratosthenes für jede Eingabe $n\geq m$ ausgegeben.
\end{center}
Wir zeigen die beiden Richtungen.

$\Rightarrow$:

Da $m$ eine Primzahl ist, besitzt sie keine echten Teiler. Damit ist sie nicht das Vielfache einer Zahl aus der Menge $\left\{2,3,4, ..., m-1\right\}$. Deshalb hat das Sieb die Zahl $m$ bis zum Erreichen von $m$ nicht durchgestrichen und somit wird $m$ vom Sieb ausgegeben.
$\Leftarrow$:

Sei $n\geq m$ die Eingabe für den Algorithmus \textit{Sieb des Eratosthenes}. Nach Konstruktion des Algorithmus, wird $m$ genau dann ausgegeben, wenn $m$ kein Vielfaches einer Zahl aus der Menge $\left\{2,3,4, ..., m-1\right\}$ ist. Damit besitzt $m$ aber keinen echten Teiler und somit ist $m$ eine Primzahl.
\item Für die Eingabe $n$ liefert das Sieb alle Primzahlen $\leq n$. Um $\text{Prim}(n)$ zu bestimmen, müssen diese Primzahlen nur noch gezählt werden.
\end{enumerate}
\end{Answer}

\clearpage

Alle nachfolgenden Python-Codes in diesem Kapitel können von Moodle heruntergeladen werden (Dateiname: {\tt prim.py}).

Die Kommentare sind (weil das so üblich ist) in englisch geschrieben.

Betrachten wir nun eine Umsetzung des Siebs von Eratosthenes in Python:
\begin{listing}[H]
\begin{minted}[linenos]{python}
def sieve(n):
    # implementation of Eratosthene's Sieve
    # siev(n) returns a list 'primes' of all prime numbers <= n

    # start with an empty list, since no prime numbers have been found yet
    primes = []

    # numpy array of type bool (either True of False) of length n+1
    # if not_crossed[i] is True, the number i has not been
    # crossed out yet and is therefore a prime number
    # not_crossed is initialized with all 'True' entries (ones):
    # not_crossed = [True, True,... , True], since no numbers have
    # been crossed out yet
    not_crossed = np.ones(n+1, dtype=bool)

    # find prime numbers in the set {2,3, ..., n}
    # remember: the numbers 0 and 1 are not prime
    for i in range(2, n+1): # check every number in {2, 3, ..., n}

        # check if the number has not been crossed out yet
        # skip the number, if it has already been crossed out
        if (not_crossed[i]):

            # the number i must be prime (it has not been crossed out)
            primes.append(i) # append i to the list of primes

            # now, we cross out every multiple of i that is <= n:
            # 2*i, 3*i, 4*i, ...
            # increment by i starting from the first multiple of i
            for m in range(2*i, n+1, i):
                not_crossed[m] = False # crossing out
    return(primes) # return the list of primes
\end{minted}
\label{listing:sieve}
\end{listing}
\clearpage

\begin{Exercise}
\begin{enumerate}[label=(\alph*)]
\item Betrachte die Python-Funktion $\mintinline{python}{def sieve(n)}$ genau und versuche zu verstehen, warum dies genau die Umsetzung des Sieb des Eratosthenes ist.
\item Lade den Python-Code {\tt prim.py} von Moodle herunter und bestimme $\text{Prim}(1000000)$.
\end{enumerate}
\end{Exercise}

\begin{Answer}

Keine Lösung notwendig.
\end{Answer}


Mit Hilfe des Sieb von Eratosthenes, ist es nun ganz einfach, eine Funktion zu definieren, welche $\text{Prim}(n)$ berechnet:
\begin{minted}[linenos]{python}
def prim(n):
    # prim(n) returns the number of primes <= n
    # implements the mathematical function Prim

    # sieve(n) returns a list of all primes <= n
    # the length (len) of this list obviously tells us
    # how many primes <= n there are
    return(len(sieve(n))) # calls our function 'sieve'
\end{minted}
Die Laufzeit des \textit{Sieb des Eratosthenes} kann an zwei Stellen deutlich verbessert werden. Die Folgende Aufgabe hilft Dir dabei, dieses Stellen zu erkennen.


\begin{Exercise}
\underline{Verbesserung 1}

Sei $n\geq 2$ eine natürliche Zahl. Zeige, dass folgende zwei Aussagen äquivalent sind:
\begin{enumerate}
    \item Die Zahl $n$ ist eine Primzahl.
    \item Die Zahl $n$ besitzt keinen echten Teiler $\leq \floor{\sqrt{n}}$.
\end{enumerate}
Dabei ist $\floor{x}$ die nächste ganze Zahl $\leq x$ (Abrundeklammer).
\end{Exercise}
\begin{Answer}

1. $\Rightarrow$ 2.

Da $n$ eine Primzahl ist, besitzt sie keinen echten Teiler und somit insbesondere auch keinen echten Teiler $\leq \floor{\sqrt{n}}$.

2. $\Rightarrow$ 1.

Sei $n$ keine Primzahl. Damit besitzt $n$ eine Faktorisierung $n = ab$ mit $a,b\in\mathbb{N}$. Wir zeigen, dass mindestens einer der Faktoren $\leq \floor{\sqrt{n}}$ ist und $n$ somit einen echten Teiler $\leq \floor{\sqrt{n}}$ hat. Falls $n$ eine Quadratzahl ist, dann besitz $n$ die Faktorisierung $n = a^2$ mit $a = \sqrt{n}$. Damit ist $a$ ein echter Teiler $\leq \floor{\sqrt{n}}$. Sei also $n$ keine Quadratzahl. Dann muss  genau einer der Beiden Faktoren $a$ oder $b$ kleiner oder gleich $\floor{\sqrt{n}}$ sein, da sonst das Produkt $ab$ grösser als $n$ wäre (im Widerspruch zu $n = ab$).
\end{Answer}



\begin{Exercise}
Wir wenden das Sieb des Eratosthenes auf die Eingabe $n>2$ an.

Angenommen, das Sieb des Eratosthenes hat gerade eine Primzahl $p > 2$ gefunden. Zeige, dass die Vielfachen $hp$ von $p$ mit $h\in \left\{2,3, ..., p-1\right\}$ bereits vom Algorithmus gestrichen wurden. Damit genügt es, die Vielfachen $\geq p^2$ zu streichen (solange sie kleiner oder gleich $n$ sind).
\end{Exercise}

\begin{Answer}
Sei $hp$ für ein $h\in \left\{2,3, ..., p-1\right\}$ ein Vielfaches von $p$. Angenommen $h$ ist eine Primzahl. Da $h<p$ wurde $h$ bereits als Primzahl vom Sieb (vor $p$) gefunden und die Vielfachen $\leq n$ von $h$ wurden gestrichen. Insbesondere also auch das Vielfache $hp$ von $h$. Falls $h$ keine Primzahl ist, dann besitzt $h$ eine Zerlegung $h = ab$ für eine Primzahl $a<h<n$. Damit wurde aber $hp = abp$ als $bp$-faches Vielfache von $a$ bereits gestrichen.
\end{Answer}

\clearpage

Mit Hilfe dieser zwei Verbesserungen, lässt sich $\text{Prim}(n)$ effizienter implementieren:
\begin{minted}[linenos]{python}
def sieve_fast(n):
    not_crossed = np.ones(n+1, dtype=bool)

    i = 2 # first prime number
    while(i**2 <= n): # Verbesserung 1
        if (not_crossed[i]):

            # Verbesserung 2
            for m in range(i**2, n+1, i):
                not_crossed[m] = False
        i = i + 1

    primes = []

    for k in range(2, n+1):
        if (not_crossed[k]):
            primes.append(k)
    return(primes)

def prim_fast(n):
    return(len(sieve_fast(n)))
\end{minted}


\begin{Exercise}
\begin{enumerate}[label=(\alph*)]
\item Bestimme $\text{Prim}(n)$ für $n = 10000000$ (zehn Millionen) mit Hilfe der Funktion \mintinline{python}{prim_fast} und vergleiche den Wert der Schätzung $n/\ln{n}$ aus dem Primzahlensatz. Wie gut ist die Schätzung?

\item Lade das Python-Programm {\tt prim\_plots.py} von Moodle herunter und führe es aus. Betrachte den generierten Plot. Was sagt der Plot aus? Wie sind die Achsen skaliert?

\item Schreibe eine Python-Funktion \mintinline{python}{def prim_approx(n)}, welche $\text{Prim}(n)$ mit Hilfe der Abschätzung, welche von dem Primzahlsatz \ref{eq:primzahlsatz} geliefert wird, approximiert.

\item Verwende die sehr schnelle Funktion $\mintinline{python}{prim_approx}$ anstelle von $\mintinline{python}{prim_fast}$, um den Wertebereich des Plots von $8\cdot 10^7$ auf $10^{10}$ zu erhöhen.
\end{enumerate}
\end{Exercise}


\begin{Answer}
\begin{enumerate}[label=(\alph*)]
\item Das Python-Programm liefert $\text{Prim}(10000000) = 664579$. Der Primzahlsatz liefert $10^7/\ln{(10^7)}\approx 620421$. Der Fehler ist etwa $7$ Prozent.

\item Beide Achsen sind logarithmiert. Der Plot zeigt, wie stark $f(n) = \text{Prim}(n)$ verglichen mit der linearen Funktion $g(n) = n$ wächst.

\item Der folgende Code berechnet die Approximation:

\begin{minted}[linenos]{python}
import numpy as np

def prim_approx(n):
    return(n/(np.log(n)))
\end{minted}
\item 
Folgende Anpassungen müssen in {\tt prim\_plots.py} durchgeführt werden:
\begin{minted}[linenos]{python}
# Plots
N = [ 10**1, 10**2, 10**3, 10**4, 10**5, 10**6, 5*10**6, 10**7,
     4*10**7, 8*10**7 , 10**10]

results = []
for n in N:
    results.append(prim_approx(n))
\end{minted}
\end{enumerate}
\end{Answer}


\begin{Exercise}
Erstelle einen zweiten Plot, welcher das Verhältnis $\text{Prim}(n) / (n / \ln{(n)})$ plottet. Was stellst Du fest? Was würde man erwarten?
\end{Exercise}

\begin{Answer}

Der entsprechende Code ist:
\begin{minted}[linenos]{python}
# quotient Prim(n)/n
ratio = []
for k in range(len(N)):
    ratio.append(results[k] / N[k])
fig, ax = plt.subplots()
ax.loglog(N, ratio, 'go-', label ='f(n) / g(n) = Prim(n) / n')
ax.set(xlabel='n', ylabel='count', title='Prim(n) / n')
ax.grid()
legend = ax.legend(loc='upper right', shadow=True, fontsize='x-large')
fig.savefig('/home/tom/Desktop/Prim_ratio.pdf')
plt.show()
\end{minted}
\end{Answer}



\subsubsection{Binäre Darstellung einer Zahl}
Dies ist das letzte mathematische Werkzeug, das wir benötigen.

Sei $m\geq 1$ eine natürliche Zahl. Wir bezeichnen mit $\text{Bin}(m)$ die kürzeste binäre Darstellung von $m$, also $\text{Nummer}(\text{Bin}(m)) = m$. Die binäre Darstellung ist offenbar genau dann am kürzesten, wenn die erste Ziffer der Darstellung eine $1$ und keine $0$ ist. Beispielsweise sind $100$ und $00100$ beides binäre Darstellungen der Zahl $4$, doch die erste ist die kürzeste.


\begin{Exercise}
Bestimme $\text{Bin}(11)$ und $\text{Bin}(52)$.
\end{Exercise}

\begin{Answer}
$\text{Bin}(11) = 1011$, $\text{Bin}(52) = 110100$
\end{Answer}



\begin{Exercise}
Schreibe eine Python-Funktion $\mintinline{python}{def Bin(n)}$, welche $\text{Bin}(n)$ für eine natürliche Zahl $n$ berechnet.
\end{Exercise}

\begin{Answer}
Lösungsvorschlag:
\begin{minted}[linenos]{python}
def Bin(n):
    x = []
    a = True
    while(a):
        r = n % 2
        x.insert(0,r)
        n = int(n/2)
        if n == 0:
            a = False
    return(x)
\end{minted}
\end{Answer}

\begin{Exercise}
Sei $m \neq 0$ eine natürliche Zahl. Beweise, dass die Länge $|\text{Bin}(m)|$ von $\text{Bin}(m)$ durch $\ceil{\log_2(m+1)}$ gegeben ist.
\end{Exercise}


\begin{Answer}
Sei $k = |\text{Bin}(m)|$. Wir wollen $k$ ermitteln. Offensichtlich gilt für $m$ die Sandwich-Beziehung
\begin{align*}
    2^k-1 \geq m > 2^{k-1} - 1 \Rightarrow \\
    2^k \geq m+1 > 2^{k-1} \Rightarrow \\
    k \geq \log_2(m+1) > k-1.
\end{align*}
Die Beziehung $k \geq \log_2(m+1) > k-1$ sagt aus, dass $\log_2(m+1)\in (k-1,k]$ und damit folgt:
\begin{align*}
    k = \ceil{\log_2(m+1)}.
\end{align*}
\end{Answer}




\clearpage
\subsection{Ein randomisiertes Kommunikationsprotokoll}
Nun sind wir in der Lage das randomisierte Kommunikationsprotokoll zu beschreiben. Zur Erinnerung: $R_1$ ist ein Rechner, der sich an der ETH in Zürich befindet (mit Datensatz $x$) und $R_2$ ist ein Rechner am MIT in Boston, USA (mit Datensatz $y$). Wir wollen prüfen, ob $x = y$ gilt (ob die Datensätze identisch sind).

Das Protokoll $R = (R_1,R_2)$ für die beiden Rechner arbeitet wie folgt.
\begin{mdframed}[backgroundcolor=green!20,roundcorner=10pt,userdefinedwidth=16cm,align=center]
\textbf{Ausgangssituation}:

$R_1$ hat $n$ Bits $x=x_1...x_n$, $R_2$ hat $n$ Bits $y = y_1...y_n$.

\textbf{Phase 1}:

$R_1$ wählt zufällig mit einer uniformen Wahrscheinlichkeitsverteilung $p$ als eine der $\text{Prim}(n^2)\approx n^2/\ln(n^2)$ (Primzahlsatz \ref{eq:quadrat}) Primzahlen kleiner gleich $n^2$.


\textbf{Phase 2}:

$R_1$ berechnet die Zahl $s = \text{Nummer}(x)\; \text{mod} \; p$ und schickt die binäre Darstellung von $s$ und $p$ and $R_2$


\textbf{Phase 3}:

Nach Empfang von $s$ und $p$ berechnet $R_2$ die Zahl $q = \text{Nummer}(y)\; \text{\text{mod}} \; p$.
\begin{itemize}
    \item Falls $q\neq s$, dann gibt $R_2$ die Ausgabe ``ungleich'' aus.
    \item Falls $q = s$, dann gibt $R_2$ die Ausgabe ``gleich'' aus.
\end{itemize}
\end{mdframed}



\begin{Exercise}
Erkläre, warum das Protokoll garantiert richtig entscheidet, falls die Ausgabe ``ungleich'' ist.
\end{Exercise}

\begin{Answer}
Da $s = \text{Nummer}(x)\; \text{mod} \; p$ und $q = \text{Nummer}(y)\; \text{\text{mod}} \; p$ die Reste der Modulo-Operation mit $p$ sind, ist es nicht möglich, dass $s$ und $q$ unterschiedlich sind, wenn $x$ und $y$ gleich sind.
\end{Answer}

Jetzt analysieren wir die Arbeit von $R = (R_1,R_2)$. Zuerst bestimmen wir die Komplexität gemessen als die Anzahl der Kommunikationsbits. Anschliessend analysieren wir die Zuverlässigkeit (Fehlerwahrscheinlichkeit) von $R = (R_1,R_2)$.

Die einzige Kommunikation besteht darin, dass $R_1$ die Zahlen $s$ und $p$ an $R_2$ schickt.

Weil $s\leq p < n^2$ ($p$ ist strikt kleiner als $n^2$, da $n^2$ offensichtlich keine Primzahl sein kann) gilt, ist die Länge der binären Nachricht $2\cdot \ceil*{\log_2(n^2)}\leq 4\cdot \ceil*{\log_2(n)}$. Für $n = 10^{16}$ sind dies höchstens $4\cdot 16\cdot \ceil*{\log_2(10)}=256$ Bits. Es ist also eine sehr kurze Nachricht, die man in problemlos zuverlässig übertragen kann.

Bei der Analyse der Fehlerwahrscheinlichkeit unterscheiden wir zwei Möglichkeiten bezüglich der tatsächlichen Relation zwischen den Datensätzen $x$ und $y$.
\begin{enumerate}
    \item Sei $x = y$. Dann gilt
    \begin{align*}
        \text{Nummer}(x)\; \text{mod} \; p = \text{Nummer}(y)\; \text{mod} \; p
    \end{align*}
    für alle Primzahlen $p$. Also gibt $R_2$ die Antwort ``gleich'' mit Sicherheit. In diesem Fall ist die Fehlerwahrscheinlichkeit $0$.
    \item Sei $x\neq y$. Wir bekommen eine falsche Antwort ``gleich'' nur dann, wenn $R_2$ eine zufällige Primzahl $p$ gewählt hat, die die Eigenschaft hat, dass
    \begin{align*}
        z = \text{Nummer}(x)\; \text{mod} \; p = \text{Nummer}(y)\; \text{mod} \; p
    \end{align*}
    gilt. In anderer Form geschrieben:
    \begin{align*}
        \text{Nummer}(x) = x'\cdot p + z
    \end{align*}
    und
    \begin{align*}
        \text{Nummer}(y) = y'\cdot p + z
    \end{align*}
    für irgendwelche natürlichen Zahlen $x'$ und $y'$.
    
    Daraus folgt, dass
    \begin{align*}
        \text{Nummer}(x)-\text{Nummer}(y) = x'\cdot p -  y'\cdot p = (x'-y')\cdot p,
    \end{align*}
    also dass $p$ die Zahl $|\text{Nummer}(x)-\text{Nummer}(y)|$ teilt.
    
    Also gibt unser Protokoll eine falsche Antwort nur, wenn die gewählte Primzahl $p$ die Zahl $|\text{Nummer}(x)-\text{Nummer}(y)|$ teilt. Wir wollen uns nun überlegen, wie wahrscheinlich dieser Fall ist.
    
    Wir wissen, dass $p$ aus $\text{Prim}(n^2)$ Primzahlen aus $\left\{2,3,...,n^2\right\}$ mit uniformer Wahrscheinlichkeitsverteilung gewählt wurde. Es ist also hilfreich festzustellen, wie viele dieser $\text{Prim}(n^2)\approx n^2/\ln(n^2)$ Primzahlen die Zahl $|\text{Nummer}(x)-\text{Nummer}(y)|$ teilen können.
    
    Da die Länge von $x$ und $y$ gleich $n$ ist, gilt
    \begin{align*}
        w = |\text{Nummer}(x)-\text{Nummer}(y)| < 2^n.
    \end{align*}
    Sei $w = p_1^{i_1}p_2^{i_2}...p_k^{i_k}$, wobei $p_1<p_2<...<p_k$ Primzahlen  und $i_1, i_2,...,i_k$ positive ganze Zahlen sind. Wir wissen, dass jede Zahl eine solche eindeutige Faktorisierung besitzt. Unser Ziel ist zu beweisen, dass $k\leq n-1$ ist.
    
    Wir beweisen diese Aussage indirekt (durch Widerspruch).
    
    Angenommen die Aussage wäre falsch, also $k\geq n$. Dann ist für genügend grosse $n$
    \begin{align*}
        w = w = p_1^{i_1}p_2^{i_2}...p_k^{i_k} \geq p_1p_2...p_n > 1\cdot 2\cdot 3\cdot ... \cdot n = n! > 2^n.
    \end{align*}
    Dies ist aber ein Widerspruch zur Tatsache, dass $w < 2^n$. Damit ist unsere Annahme falsch und $w$ kann höchstens $n-1$ unterschiedliche Primfaktoren haben. Weil jede Primzahl aus $\left\{2,3,...,n^2\right\}$ die gleiche Wahrscheinlichkeit hat, (zufällig) gewählt zu werden, ist die Wahrscheinlichkeit, ein $p$ zu wählen, das $w$ teilt, höchstens
\begin{align*}
    \frac{n-1}{\text{Prim}(n^2)}\leq \frac{n-1}{n^2/\ln(n^2)}\leq\frac{\ln(n^2)}{n}
\end{align*}
    für genügend grosse $n$.
    
    Also ist die Fehlerwahrscheinlichkeit des Kommunikationsprotokolls höchstens $\ln(n^2)/n$, was für $n=10^{16}$ höchstens $0.736827\cdot 10^{-14}$ ist.
\end{enumerate}
Sollte uns diese Fehlerwahrscheinlichkeit noch zu hoch sein, kann das Protokoll z.B. 10 mal wiederholt werden (mit 10 zufälligen Primzahlen) und man erhält dann eine Fehlerwahrscheinlichkeit von höchstens
    \begin{align*}
        \left(\frac{n-1}{\text{Prim}(n^2)}\right)^{10}\leq \left(\frac{\ln(n^2)}{n}\right)^{10} = \frac{2^{10}\cdot (\ln(n))^{10}}{n^{10}}
    \end{align*}
    für genügend grösse $n$. Für $n = 10^{16}$ ist dies höchstens $0.472\cdot 10^{-141}$. Eine Fehlerwahrscheinlichkeit, die man gerne in Kauf nimmt.
    
\begin{Exercise}
Verwende deine Python-Funktionen $\mintinline{python}{sieve_fast}$, $\mintinline{python}{Bin}$ und $\mintinline{python}{Nummer}$ um das Kommunikationsprotokoll in Python zu implementieren.
\end{Exercise}
\begin{Answer}

Die Funktionen \mintinline{python}{sieve_fast}, \mintinline{python}{Nummer} und \mintinline{python}{Bin} haben wir schon geschrieben. Diese können wir nun hier verwenden. Die Funktion \mintinline{python}{is_equal} wird nicht benötigt, kann aber verwendet werden um (deterministisch) zu prüfen, ob zwei Datensätze identisch sind oder nicht.
\begin{minted}[linenos]{python}
import numpy as np

def sieve_fast(n):
    not_crossed = np.ones(n+1, dtype=bool)

    i = 2
    while(i**2 <= n):
        if (not_crossed[i]):
            for m in range(i**2, n+1, i):
                not_crossed[m] = False
        i = i + 1

    primes = []

    for k in range(2, n+1):
        if (not_crossed[k]):
            primes.append(k)
    return(primes)

def Nummer(x):
    x = list(str(x))
    n = len(x)
    dec = 0
    for i, xi in enumerate(x):
        dec = dec + int(xi) * 2**(n-(i+1))
    return(dec)



def Bin(n):
    x = []
    a = True
    while(a):
        r = n % 2
        x.insert(0,r)
        n = int(n/2)
        if n == 0:
            a = False
    return(x)

def is_equal(x,y):
    nx = len(x)
    ny = len(y)

    if nx != ny:
        return(False)
    for i in range(nx):
        if x[i] != y[i]:
            return(False)
    return(True)
\end{minted}

Das Protokoll lässt sich nun recht einfach wie folgt umsetzen:
\begin{minted}[linenos]{python}
def pick_random_prime(n):
    primes = sieve_fast(n)
    p = np.random.choice(primes,1)[0]
    return(p)

def R1(x):
    n = len(str(x))
    p = pick_random_prime(n**2)
    s = Nummer(x) % p
    return([s,p])

def protokoll(x,y):
    sp = R1(x)
    q = Nummer(y) % sp[1]
    if sp[0] == q:
        return(True)
    else:
        return(False)

x = '100110101'
y = '100110101'

print(protokoll(x,y))
\end{minted}
\end{Answer}
